\chapter{Anoscopy}

\section*{Overview}
Anoscopy is an examination of the anus and distal rectum using a short, rigid, lighted instrument called an anoscope. It provides direct visualization of the anal canal and anorectal junction. The exact approach, setting, and preparation depend on the indication, anatomy, and the patient's health.

\section*{Alternative Names}
Anal examination, Anal canal inspection, Proctoscopy

\section*{Principle}
The anoscope is a short, rigid tube with an obturator and a light source that is gently inserted through the anal sphincter. It holds the anal canal open, allowing direct inspection of the mucosa, hemorrhoidal cushions, and anorectal junction. Biopsy forceps and therapeutic instruments can be passed through the scope to sample or treat lesions.

\section*{History}
Direct speculum examination of the rectum dates to the 19th century. Modern anoscopes with integrated lighting became standard in the mid-20th century.

\section*{Why It Is Done}
Anoscopy is ordered to evaluate anorectal complaints such as rectal bleeding, anal pain, itching, discharge, or a palpable mass. It is used to diagnose internal and external hemorrhoids, anal fissures, polyps, abscesses, fistulas, and anorectal tumors, and to obtain biopsies of suspicious lesions.

\section*{Is This a Primary or Secondary Test or Procedure}
Primary test for evaluating symptoms and signs localized to the distal anal canal.

\section*{How It Is Performed}
The patient lies on the side with knees drawn toward the chest, and the perianal skin is inspected first. The lubricated anoscope is gently inserted through the anus and advanced into the rectum, then withdrawn slowly while the examiner inspects the mucosa. Biopsies or therapeutic interventions can be performed through the scope.

\section*{Preparation}
No special preparation is required; the patient may be asked to empty the bowel beforehand. A small enema may be used to clear the anal canal. Sedation is usually not needed.

\section*{Contraindications and Precautions}
Use caution in patients with recent anorectal surgery, severe anal stenosis, or an acutely thrombosed, very painful hemorrhoid. Acute anal fissure with severe spasm may make instrumentation difficult. The examination is generally avoided when anorectal infection is suspected to require imaging first.

\section*{Risks}
Minor bleeding, transient discomfort, or mild pain after the examination. Injury to the anal mucosa and infection are uncommon.

\section*{Aftercare and Recovery}
Most patients resume normal activities immediately after the examination. Minor spotting or discomfort may persist for a few hours. Persistent bleeding, severe pain, or fever should be reported promptly.

\section*{Outcome and Follow-up}
Findings are interpreted by direct visual inspection, with suspicious lesions sampled for histologic confirmation. Results help distinguish conditions such as hemorrhoids, fissures, and malignancy.

\section*{Expected Results}
Normal appearance of the anal mucosa without lesions, inflammation, masses, or abnormal discharge.

\section*{Follow-up}
Colonoscopy is recommended if proximal colonic disease is suspected or screening is indicated. Repeat anoscopy may be used to monitor healing or the response of known lesions to treatment.

\section*{When to Seek Medical Care}
Follow the procedure team's specific instructions. Seek urgent medical assessment for severe or worsening pain, uncontrolled bleeding, fever or signs of infection, trouble breathing, chest pain, fainting, new weakness or confusion, inability to urinate, or any rapidly worsening symptom. The appropriate warning signs depend on the procedure and the patient's medical condition.

\section*{References}
\begin{enumerate}
  \item MedlinePlus. Anoscopy. U.S. National Library of Medicine; 2025.
  \item Current Medical Diagnosis and Treatment. McGraw-Hill; 2025.
\end{enumerate}

\clearpage
